\begin{abstract}
Hypervisors are traditionally inseparable from a particular host operating
system (OS) or bare-metal environment: KVM is woven into Linux,
static-partitioning hypervisors such as Bao and Jailhouse assume specific boot
and device-management models, and even research designs like DuVisor bind the
virtualization layer to a single kernel substrate. This coupling forces every
new or emerging OS that needs virtualization to either reimplement a hypervisor
from scratch or deeply restructure its kernel, which fragments engineering
effort and impedes reuse of mature virtualization logic. We argue that this
coupling is a software-architecture problem rather than a fundamental
virtualization requirement: a hypervisor depends on its substrate OS through a
small but semantically rich set of services. We present \axvisor{}, a hypervisor
restructured around an explicit \emph{substrate interface contract} that
decouples the virtualization core from OS-specific services. The same
\axvisor{} core then runs on three substantially different
substrates---\arceos{}, \asterinas{}, and Linux---each contributing only a thin
adaptation layer. To decouple the guest-facing side, \axvisor{} additionally
exposes a KVM-compatible ABI, allowing unmodified guests and toolchains built
for the de facto KVM interface to run on top of any substrate. We evaluate the
design by booting a bare-metal Linux guest on both the \asterinas{} and Linux
substrates, quantifying portability cost through interface size and code-reuse
metrics, and measuring virtualization overhead. Our results show that
substrate-OS portability is achievable with a narrow contract and modest
per-substrate adaptation, turning the hypervisor from an OS-specific subsystem
into a reusable, retargetable virtualization component.
\end{abstract}
